\chapter{Reprodução dos resultados}
\label{ap:reproducao}

Todos os cálculos estão no repositório do projeto. A
Tabela~\ref{tab:ap-reproducao} relaciona cada resultado aos programas
e arquivos que o produzem. Os caminhos são relativos à raiz do
repositório.

\begin{table}[htbp]
\centering
\caption{Resultados e programas correspondentes.}
\label{tab:ap-reproducao}
\small
\begin{tabular}{p{0.36\textwidth}p{0.56\textwidth}}
\toprule
Resultado & Programas e dados\\
\midrule
Matriz de marés, postos, vínculo escalar, expressões do TG
(Cap.~\ref{cap:polarizacoes}) &
\texttt{src/polarizacoes/foundations.py}; \texttt{scripts/reproduzir\_tg.py}
(19 testes simbólicos); \texttt{tests/test\_polarizacoes.py},
\texttt{tests/test\_normalizacao.py}\\[1mm]
Correlações tensoriais com termos dos pulsares; validação por duas rotas
(Cap.~\ref{cap:pta}) &
\texttt{src/pta/orf.py}, \texttt{src/pta/response.py};
\texttt{scripts/validar\_orf.py}, \texttt{scripts/benchmark\_orf.py};
\texttt{results/C05/}\\[1mm]
Figuras desta versão e conferência independente das correlações
apenas-Terra (tensorial e escalar) &
\texttt{dissertacao\_v2/scripts/figuras.py}\\[1mm]
Campanha tensorial I (ruído livre) &
\texttt{scripts/run\_calibration\_campaign.py},
\texttt{scripts/sintetizar\_calibracao.py}; \texttt{results/C07/}\\[1mm]
Campanha tensorial II e contrastes condicionais &
\texttt{scripts/comparar\_populacao\_c13.py},
\texttt{scripts/contrastes\_condicionais\_c13.py};
\texttt{results/C11/statistical\_audit.json},
\texttt{results/C13/conditional\_contrasts/}\\[1mm]
Campanha escalar, falsas detecções e Fisher &
\texttt{scripts/gerar\_populacao\_escalar\_c10.py},
\texttt{scripts/sintetizar\_campanha\_escalar\_c10.py},
\texttt{scripts/calcular\_fisher\_escalar\_c10.py};
\texttt{results/C10/population\_synthesis/}\\[1mm]
Experimento de ajuste de temporização &
\texttt{scripts/experimento\_timing\_c13.py};
\texttt{results/C13/timing\_projection/}\\
\bottomrule
\end{tabular}
\end{table}

\section{Controle do erro numérico}

As posteriores foram calculadas por quadratura determinística na massa
e, na campanha escalar, também em $\varepsilon$. Não se usou amostragem
por Monte Carlo. Para cada realização, a posterior foi avaliada com duas
resoluções, e a diferença entre elas foi propagada como um intervalo
numérico em cada quantidade derivada, como quantis, coberturas e fatores
de Bayes. Quando o intervalo numérico não permitia decidir um teste, o
resultado foi registrado como indeterminado, e a realização não foi
descartada. As contagens citadas no Capítulo~\ref{cap:simulacoes} são
os valores nominais. Em todos os casos citados, a conclusão qualitativa
não depende da resolução numérica. As exceções, como as comparações
marginalizadas da aproximação de frequência de referência, foram
declaradas inconclusivas no texto.

As correlações angulares foram aceitas apenas após refinamento
independente da quadratura angular e do corte multipolar, com tolerância
absoluta de $10^{-5}$ e relativa de $10^{-3}$.
